% ==========================================================
% 第6章 局所化と局所環
% ==========================================================
\chapter{局所化と局所環}
\label{ch:ch06}

\begin{chapterlead}
第5章では，アフィン多様体 $X$ 全体の情報が座標環 $k[X]$ に集約されることを見た．
しかし幾何学の多くの問い ―― 点 $p$ は特異点か，$p$ の近くで $X$ は何本の枝をもつか，$p$ での接空間の次元はいくつか ――
は\emph{点の近く}だけで決まる局所的な問いである．
本章では，環から「点の近くの情報」だけを取り出す代数的操作である\textbf{局所化}を導入し，
その結果得られる\textbf{局所環} $\OO_{X,p}$ を調べる．
局所化とは一言でいえば「ある元たちを可逆にする（分母として許す）」操作であり，
幾何学的には「それらの関数が $0$ にならない場所へ注目する」ことに対応する．

物理学でも「局所性」は根本的な原理であるが，代数幾何の「局所」とはかなり意味が異なる．
本章の最後でその違いを丁寧に論じる．
\end{chapterlead}

本章では環はすべて単位元をもつ可換環とする．$k$ は代数閉体である．

% ----------------------------------------------------------
\section{局所的な情報を取り出す}
% ----------------------------------------------------------

\subsection{局所的な幾何学の考え方}

微分幾何では，多様体 $M$ の点 $p$ における性質（接空間，曲率，特異性など）を調べるとき，
$p$ の小さな開近傍 $U$ をとって局所座標で計算する．
$p$ を含む開集合をどんどん小さくしても変わらない情報が「$p$ における局所的な情報」である．

代数幾何でも同じ発想をとりたいが，二つの障害がある．
\begin{enumerate}
  \item Zariski 位相は非常に粗く，開集合は「巨大」である．既約な $X$ の空でない開集合はすべて稠密であり，「小さな近傍」はとれない．
  \item 多項式関数（正則関数）は剛性が強く，局所的に定義された関数を大域的に延ばせるとは限らない一方，延ばせるならその延長は一意である．
\end{enumerate}
それでも「$p$ の近くで定義された関数」を全部集めてくれば，$p$ の周辺の構造を反映した環が得られる．
それを代数的に構成するのが局所化である．

\subsection{点の近傍}

$X$ をアフィン多様体，$k[X]$ をその座標環とする．$f\in k[X]$ に対し
\[
  D(f)=X\setminus\Zero_X(f)=\{p\in X\mid f(p)\neq0\}
\]
を\textbf{基本開集合}という．

\begin{myprop}{基本開集合は開基をなす}{ch06-basic-open}
$\{D(f)\}_{f\in k[X]}$ は $X$ の Zariski 位相の開基をなす．さらに $D(f)\cap D(g)=D(fg)$ である．
\end{myprop}

\begin{proof}
開集合 $U=X\setminus \Zero_X(J)$ と $p\in U$ をとる．$p\notin\Zero_X(J)$ なので $f(p)\neq0$ なる $f\in J$ があり，
$\Zero_X(J)\subset\Zero_X(f)$ より $p\in D(f)\subset U$．後半は $(fg)(p)\ne0\iff f(p)\ne0$ かつ $g(p)\ne0$ から明らか．
\end{proof}

したがって点 $p$ の近傍としては $p\in D(f)$，すなわち $f(p)\ne0$ なる $f$ に対する $D(f)$ を考えれば十分である．
$p$ の「十分小さい近傍」とは「$f(p)\neq0$ なるいくつかの $f$ の非零集合の共通部分」であり，
上の命題からそれはまた一つの $D(f)$ の形になる．

\begin{figure}[H]
\centering
\begin{tikzpicture}[scale=0.9]
  \draw[thick,blue,domain=-2.6:2.6,samples=60] plot (\x,{0.25*\x*\x*\x-0.9*\x});
  \node[blue] at (2.9,2.3) {$X$};
  \fill (0.6,{0.25*0.216-0.54}) circle (2.2pt) node[below right]{$p$};
  \foreach \a in {-2.2,-1.1,1.7}{
    \draw[thick,fill=white] (\a,{0.25*\a*\a*\a-0.9*\a}) circle (2.4pt);
  }
  \node[align=left,font=\small] at (5.6,0) {$\circ$：$f$ の零点（除かれた点）\\$D(f)=X\setminus\{\circ\}$ は\\$p$ の Zariski 開近傍};
\end{tikzpicture}
\caption{曲線上の点 $p$ の Zariski 開近傍は，有限個の点を除いたもの．どれも「大きい」．}
\label{fig:ch06-nbhd}
\end{figure}

第5章の大域正則関数の定理（定理\ref{thm:ch05-global-regular}）と同様の議論で，$D(f)$ 上の正則関数は $g/f^m$（$g\in k[X]$）の形をしていることが分かる．

\begin{myprop}{$D(f)$ 上の正則関数}{ch06-Df}
$X$ をアフィン多様体，$0\neq f\in k[X]$ とする．このとき
\[
  \OO_X(D(f))=\{\,g/f^m\mid g\in k[X],\ m\ge0\,\}\subset k(X).
\]
さらに $D(f)$ はアフィン多様体 $\{(p,t)\in X\times\A^1\mid f(p)t=1\}$ と同型である．
\end{myprop}

\begin{proof}[証明の概略]
右辺が左辺に含まれるのは明らか．逆に $h\in\OO_X(D(f))$ をとり，$J=\{a\in k[X]\mid ah\in k[X]\}$ とおく．
第5章の定理（大域正則関数）と同じ議論により $\Zero_X(J)\subset\Zero_X(f)$ が分かるので，零点定理より $f\in\sqrt J$，
すなわちある $m$ で $f^m\in J$，$h=(f^mh)/f^m$．後半は第5章の例 $\A^1\setminus\{0\}\cong\Zero(xy-1)$ と同様で，
$p\mapsto(p,1/f(p))$ とその逆写像（射影）が互いに逆の射になる．
\end{proof}

この環 $\{g/f^m\}$ が次節で定義する局所化 $k[X]_f$ である．

\subsection{局所的な関数}
\label{subsec:ch06-germ}

点 $p$ の近くで定義された関数とは，$p$ のある開近傍 $U$ 上の正則関数 $f\in\OO_X(U)$ のことである．
二つの局所関数 $(U,f)$ と $(V,g)$ は，$p$ を含むより小さな開集合 $W\subset U\cap V$ 上で一致するとき「$p$ で同じ芽をもつ」という．

\begin{mydef}{関数の芽}{ch06-germ}
$p\in X$ に対し，組 $(U,f)$（$U\ni p$ は開集合，$f\in\OO_X(U)$）の全体に
\[
  (U,f)\sim(V,g)\iff \text{ある開集合 } W\ (p\in W\subset U\cap V) \text{ で } f|_W=g|_W
\]
という同値関係を入れる．同値類を $p$ における正則関数の\textbf{芽}といい，芽全体を $\OO_{X,p}$ と書く．
\end{mydef}

芽どうしの和と積は，共通部分に制限して計算すれば well-defined であり，$\OO_{X,p}$ は $k$ 代数になる．
芽 $[U,f]$ の $p$ での値 $f(p)$ も well-defined である．

既約な $X$ では一致の定理（第5章）により，空でない開集合上で一致する正則関数は定義域全体で一致する．
したがって $(U,f)\sim(V,g)\iff f|_{U\cap V}=g|_{U\cap V}$ であり，芽は関数体 $k(X)$ の元とみなせる：
\[
  \OO_{X,p}=\{\,g/h\in k(X)\mid g,h\in k[X],\ h(p)\neq0\,\}.
\]
すなわち $\OO_{X,p}$ は「$p$ で分母が消えない有理関数」全体である．
滑らかな関数の芽（$C^\infty$ 芽）とは異なり，代数的な芽は芽から関数が（定義域の上で）完全に復元されることに注意する．

\subsection{局所的な情報と大域的な情報}

以上の考察により，座標環 $k[X]$ から出発して次の三段階の環が得られる：
\[
  k[X]=\OO_X(X)\ \subset\ \OO_X(D(f))=k[X]_f\ \subset\ \OO_{X,p}\ \subset\ k(X)\qquad(f(p)\neq0).
\]
左ほど大域的，右ほど局所的であり，右に行くほど「割ってよい元」が増える．
最も右の関数体 $k(X)$ は「どこか空でない開集合上で定義された関数」全体で，これは $X$ の生成点（第12章）での局所環にあたる．

大域的な情報と局所的な情報の関係は単純ではない．例えば
\begin{itemize}
  \item すべての点 $p$ で $\OO_{X,p}$ が「良い」（正則局所環，第7章）ことと，$X$ が非特異であることは同値である（局所的な性質）．
  \item 一方 $\OO_{X,p}$ をすべての点で知っていても，それらがどう貼り合わさっているかという大域的な情報は別に必要である．
        実際，任意の二つの非特異曲線 $X,Y$ の局所環 $\OO_{X,p}$，$\OO_{Y,q}$ は（完備化すれば）どれも $k[[t]]$ と同型になってしまう．
\end{itemize}
局所的なデータを貼り合わせて大域的な対象を記述する枠組みが層であり，第11章で導入する．

\begin{physbox}[局所座標と場]
一般相対論では時空の点 $x$ の近傍で局所慣性系をとり，そこでの物理（等価原理）を議論する．
これは「点の近傍の芽」で物理を記述する発想である．
ただし代数幾何の Zariski 近傍は巨大であり，局所慣性系のような「小さな近傍」ではない．
代数幾何で真に「無限小近傍」に対応するのは，局所環をさらに完備化した $\widehat{\OO}_{X,p}$（形式的冪級数）であり，
これは物理の Taylor 展開・摂動展開に近い（6.4 節，6.5 節）．
\end{physbox}

\exprob{prob:ch06-1-1}
$X=\A^1$，$p=0$ とする．芽 $[\A^1\setminus\{1\},\ 1/(1-x)]\in\OO_{\A^1,0}$ と，形式的冪級数 $\sum_{n\ge0}x^n$ の関係を説明せよ．
また $\OO_{\A^1,0}$ の元で，$\A^1$ 全体の多項式関数にはならないものの例を挙げよ．

% ----------------------------------------------------------
\section{局所化}
% ----------------------------------------------------------

\subsection{局所化の動機}

整数環 $\Z$ から有理数体 $\Q$ を作る操作は，「$0$ 以外のすべての整数で割ることを許す」ものであった．
その途中として，例えば「$2$ の冪でのみ割ることを許す」環 $\Z[1/2]=\{a/2^m\}$ や，
「素数 $p$ と互いに素な数で割ることを許す」環 $\Z_{(p)}=\{a/b\mid p\nmid b\}$ を考えることができる．

幾何学的には，前節で見たように
\begin{itemize}
  \item $D(f)$ 上の関数 $=$ 「$f$ の冪で割ることを許した」関数，
  \item 点 $p$ における芽 $=$ 「$p$ で消えない関数すべてで割ることを許した」関数
\end{itemize}
であった．いずれも「指定した集合 $S$ の元を分母として許す」という操作であり，これを一般の可換環で定式化するのが局所化である．

\subsection{積閉集合}

分母として許す集合は，積について閉じていなければならない（$1/s$ と $1/t$ があれば $1/(st)$ もあるべき）．

\begin{mydef}{積閉集合}{ch06-mult-set}
環 $R$ の部分集合 $S$ が\textbf{積閉集合}であるとは，$1\in S$ であり，$s,t\in S\Rightarrow st\in S$ を満たすことをいう．
\end{mydef}

\begin{myex}{積閉集合の基本例}{ch06-mult-set-ex}
\begin{enumerate}
  \item $f\in R$ に対し $S_f=\{1,f,f^2,\dots\}$．
  \item 素イデアル $\p$ に対し $S=R\setminus\p$．$s,t\notin\p\Rightarrow st\notin\p$ は素イデアルの定義そのものである．
  \item 整域 $R$ に対し $S=R\setminus\{0\}$（(2) で $\p=(0)$ の場合）．
  \item $R$ の非零因子全体．
  \item $1+I=\{1+a\mid a\in I\}$（$I$ はイデアル）．
\end{enumerate}
\end{myex}

(2) により，「素イデアルの補集合」と「積閉集合」は密接に結びついている．逆に次が成り立つ．

\begin{mylem}{積閉集合と交わらない極大なイデアル}{ch06-max-disjoint}
$S$ を $0\notin S$ なる積閉集合とする．$S$ と交わらないイデアルの中で（包含について）極大なものは素イデアルである．
\end{mylem}

\begin{proof}
極大なもの $\p$ は Zorn の補題で存在する．$a,b\notin\p$ とすると，$\p+(a)$ と $\p+(b)$ は極大性から $S$ と交わり，
$s=p_1+ra$，$t=p_2+r'b$（$p_i\in\p$）と書ける．$st\in S$ かつ $st\in\p+(ab)$ なので，$ab\in\p$ なら $st\in\p$ となり矛盾．よって $ab\notin\p$．
\end{proof}

\subsection{局所化の定義}

$S\subset R$ を積閉集合とする．分数 $r/s$（$r\in R,\ s\in S$）を形式的に作りたい．
整域なら $r/s=r'/s'\iff rs'=r's$ と定義すればよいが，一般の環（零因子をもつ環）ではこの定義は推移律を満たさない．
そこで次のように定義する．

\begin{mydef}{局所化}{ch06-localization}
$R\times S$ 上の関係を
\[
  (r,s)\sim(r',s')\iff \text{ある } u\in S \text{ が存在して } u(rs'-r's)=0
\]
で定める．これは同値関係であり，$(r,s)$ の同値類を $r/s$ と書く．同値類全体 $S^{-1}R$ に
\[
  \frac rs+\frac{r'}{s'}=\frac{rs'+r's}{ss'},\qquad \frac rs\cdot\frac{r'}{s'}=\frac{rr'}{ss'}
\]
で和と積を入れた環を $R$ の $S$ による\textbf{局所化}（分数環）という．
自然な準同型 $\iota:R\to S^{-1}R$，$r\mapsto r/1$ がある．
\end{mydef}

\paragraph{推移律と $u$ の必要性．}
反射律・対称律は明らかである．推移律を確かめよう．$(r,s)\sim(r',s')$，$(r',s')\sim(r'',s'')$ とし，
$u(rs'-r's)=0$，$v(r's''-r''s')=0$（$u,v\in S$）とする．第一式に $vs''$，第二式に $us$ を掛けて足すと
\[
  uvs''rs'-uvs''r's+uvsr's''-uvsr''s'=uvs'(rs''-r''s)=0 .
\]
$uvs'\in S$ なので $(r,s)\sim(r'',s'')$．
ここで $r's$ の項が打ち消しあったあとに残る係数 $s'$ を $S$ の元として「吸収」できることが本質であった．

もし $u$ なしで $rs'=r's$ と定義すると，上の計算は $s'(rs''-r''s)=0$ までしか言えず，$s'$ が零因子のとき $rs''=r''s$ が従わない．
具体例：$R=k[x,y]/(xy)$，$S=\{1,x,x^2,\dots\}$ とすると，
$(y,1)$ と $(0,x)$ は $y\cdot x=0\cdot 1$ を満たし，$(0,x)$ と $(0,1)$ も $0\cdot1=0\cdot x$ を満たすが，
$(y,1)$ と $(0,1)$ は $y\cdot1\ne0$．つまり「$u$ なし」の関係は推移的でない．
正しい定義では $x\cdot(y\cdot1-0\cdot1)=xy=0$ なので $y/1=0$ となる：
\emph{$x$ を可逆にすると，$x$ に零化される元は $0$ になる．}
幾何学的には，$D(x)$ 上では $y$ 軸成分が除かれるので $y$ は恒等的に $0$ になる，ということである．

和・積が well-defined であること，$S^{-1}R$ が環の公理を満たすことは同様の計算で確かめられる（演習）．
$0$ 元は $0/1$，単位元は $1/1$ であり，$s\in S$ に対し $s/1$ は逆元 $1/s$ をもつ．

\subsection{局所化された環}

基本的な性質をまとめる．

\begin{myprop}{局所化の基本性質}{ch06-loc-basic}
$S\subset R$ を積閉集合，$\iota:R\to S^{-1}R$ を自然な準同型とする．
\begin{enumerate}
  \item $\iota(s)$（$s\in S$）は $S^{-1}R$ の単元である．
  \item $\Ker\iota=\{r\in R\mid \text{ある } s\in S \text{ で } sr=0\}$．とくに $R$ が整域で $0\notin S$ なら $\iota$ は単射で，$S^{-1}R\subset\Frac R$．
  \item $S^{-1}R=0\iff 0\in S$．
  \item $S^{-1}R$ の任意の元は $\iota(r)\iota(s)^{-1}$ と書ける．
\end{enumerate}
\end{myprop}

\begin{proof}
(2) $r/1=0/1\iff$ ある $u\in S$ で $u r=0$．(3) $0\in S$ なら $1/1=0/1$（$0\cdot(1-0)=0$），逆に $1/1=0/1$ ならある $u\in S$ で $u=0$．
(1)(4) は定義から明らか．
\end{proof}

局所化のもっとも重要な性質は，イデアル（特に素イデアル）が簡明に振る舞うことである．

\begin{myprop}{局所化と素イデアル}{ch06-loc-prime}
写像 $\q\mapsto\iota^{-1}(\q)$ は全単射
\[
  \{S^{-1}R\text{ の素イデアル}\}\overset{1:1}{\longleftrightarrow}\{R\text{ の素イデアル } \p \text{ で } \p\cap S=\emptyset\}
\]
を与え，逆写像は $\p\mapsto S^{-1}\p=\{a/s\mid a\in\p,\ s\in S\}$ である．この対応は包含関係を保つ．
\end{myprop}

\begin{proof}[証明の概略]
$\q$ が素なら $\iota^{-1}(\q)$ も素であり，$S$ の元は単元に写るので $\q$ に入らず $\iota^{-1}(\q)\cap S=\emptyset$．
逆に $\p\cap S=\emptyset$ なる素イデアル $\p$ について，$S^{-1}\p$ が素イデアルであること，$\iota^{-1}(S^{-1}\p)=\p$ を示す．
$a/s\in S^{-1}\p$（$a\in\p$）と $r/1=a/s$ なら $u(rs-a)=0$，$urs=ua\in\p$ で $u,s\notin\p$ より $r\in\p$．
これから $\iota^{-1}(S^{-1}\p)=\p$ と $1\notin S^{-1}\p$ が従う．素であることも同様．
$S^{-1}(\iota^{-1}\q)=\q$ は (4) から従う．
\end{proof}

幾何学的には，$\Spec$（第12章）の言葉で「$S^{-1}R$ のスペクトルは $R$ のスペクトルのうち $S$ の元が消えない部分」であり，
局所化が開集合（やそれよりさらに小さい部分）への制限に対応することを示している．
また局所化は完全関手である（$S^{-1}$ は短完全列を保つ）という重要な性質をもつが，本ノートでは必要になった箇所で触れる．

\subsection{局所化の普遍性}

\begin{mythm}{局所化の普遍性}{ch06-universal}
$S\subset R$ を積閉集合とする．環準同型 $\phi:R\to T$ で，すべての $s\in S$ について $\phi(s)$ が $T$ の単元となるものに対し，
環準同型 $\tilde\phi:S^{-1}R\to T$ で $\tilde\phi\circ\iota=\phi$ を満たすものがただ一つ存在する．
\[
\begin{tikzcd}[column sep=large,row sep=large]
  R \arrow[r,"\iota"] \arrow[rd,"\phi"'] & S^{-1}R \arrow[d,dashed,"\exists!\,\tilde\phi"]\\
  & T
\end{tikzcd}
\]
\end{mythm}

\begin{proof}
一意性：$\tilde\phi(r/s)=\tilde\phi(\iota(r)\iota(s)^{-1})=\phi(r)\phi(s)^{-1}$ と決まってしまう．
存在：$\tilde\phi(r/s)=\phi(r)\phi(s)^{-1}$ と定義する．well-definedness：$u(rs'-r's)=0$ なら
$\phi(u)(\phi(r)\phi(s')-\phi(r')\phi(s))=0$ で $\phi(u)$ は単元なので $\phi(r)\phi(s)^{-1}=\phi(r')\phi(s')^{-1}$．準同型であることは容易．
\end{proof}

すなわち $S^{-1}R$ は「$S$ の元を可逆にする最も普遍的な $R$ 代数」である．
普遍性から $S^{-1}R$ は同型を除いて一意に定まるので，具体的な構成によらない議論ができる．
例えば次の系はこの普遍性から直ちに従う．

\begin{mycor}{$R_f$ の表示}{ch06-Rf}
$f\in R$ に対し $R_f:=S_f^{-1}R$（$S_f=\{f^m\}$）とおくと，$R_f\cong R[t]/(ft-1)$．
\end{mycor}

\begin{proof}
$R[t]/(ft-1)$ では $f$ は可逆（逆元 $t$）であり，$R\to T$ で $f$ が可逆に写るなら $t\mapsto\phi(f)^{-1}$ によって一意に $R[t]/(ft-1)\to T$ が延びる．
両者は同じ普遍性をもつので同型である．
\end{proof}

これは第5章の例 $\A^1\setminus\{0\}\cong\Zero(xy-1)$ の代数的な言い換えであり，$D(f)\cong\Zero(f t-1)\subset X\times\A^1$ に対応する．

\subsection{局所化の具体例}

\begin{myex}{局所化の例}{ch06-loc-ex}
\begin{enumerate}
  \item $R$ 整域，$S=R\setminus\{0\}$：$S^{-1}R=\Frac R$．例えば $\Z\to\Q$，$k[X]\to k(X)$．
  \item $\Z[1/2]=S_2^{-1}\Z=\{a/2^m\}$（$S_2=\{1,2,4,\dots\}$，2進分数．2進整数環や $\Z/2\Z$ と混同しないこと）．物理で浮動小数点数が表せるのはこの環の元である．
  \item $k[x]_x=k[x,x^{-1}]$：$\A^1\setminus\{0\}$ の座標環（Laurent 多項式環）．
  \item $R=k[x,y]/(xy)$，$f=x$：$R_x\cong k[x,x^{-1}]$．$y=xy/x=0$ となるので $y$ 軸成分が消え，$D(x)=\{x\ne0\}$ は $x$ 軸から原点を除いたものになる．
  \item $R=k[x]/(x^2)$，$f=x$：$x$ は冪零なので $1/1=x^2/x^2=0$，$R_x=0$．幾何学的に $D(x)=\emptyset$ である．
        一般に $R_f=0\iff f$ は冪零．
  \item $R=\Z/6\Z$，$S=\{1,3\}$（$3^2=3$ なので積閉）：$2\cdot3=0$ より $2/1=0$ となる．
        $\Z/6\Z\cong\Z/2\Z\times\Z/3\Z$ のもとで $3\mapsto(1,0)$ なので，$3$ を可逆にすると $\Z/3\Z$ 成分が消え，$S^{-1}R\cong\Z/2\Z$．
\end{enumerate}
\end{myex}

例 (6) は「可逆にすると零因子の相手が消える」現象の典型である．$\Z/6\Z$ の $\Spec$ は2点 $\{(2),(3)\}$ からなり，
$3$ を可逆にすると $3$ を含まない素イデアル $(2)$ の点だけが残り $\Z/2\Z$ になる．

\exprob{prob:ch06-2-1}
$R=k[x,y]/(xy)$，$S=\{1,x+y,(x+y)^2,\dots\}$ とする．$S^{-1}R$ を計算し，$D(x+y)\subset\Zero(xy)$ の幾何学的な形と比較せよ．

\exprob{prob:ch06-2-2}
局所化 $S^{-1}R$ において，和 $\frac rs+\frac{r'}{s'}=\frac{rs'+r's}{ss'}$ が代表元のとり方によらないことを示せ．

% ----------------------------------------------------------
\section{素イデアルによる局所化}
% ----------------------------------------------------------

\subsection{素イデアルに対する局所化}

点 $p\in X$ での関数の芽は「$p$ で消えない関数で割ることを許す」ものだった．
$p$ で消えない関数全体は $k[X]\setminus\m_p$ である．一般化して次のように定義する．

\begin{mydef}{素イデアルによる局所化}{ch06-Rp}
$\p\subset R$ を素イデアルとし，$S=R\setminus\p$ とおく．$S^{-1}R$ を $R_\p$ と書き，$R$ の $\p$ における\textbf{局所化}という．
また $R$ 加群 $M$ についても同様に $M_\p=S^{-1}M$ と書く．
\end{mydef}

記号の混乱に注意：$R_f$ は「$f$ を可逆にする」，$R_\p$ は「$\p$ の\emph{外}を可逆にする」．
$R_\p$ では $\p$ の元以外はすべて可逆になる．

\subsection{$R_\p$ の構成}

具体的に $R_\p=\{a/s\mid a\in R,\ s\notin\p\}$ である．
$R_\p$ の中で $\p R_\p=\{a/s\mid a\in\p,\ s\notin\p\}$ を考える．

\begin{myprop}{$R_\p$ はただ一つの極大イデアルをもつ}{ch06-Rp-local}
$\p R_\p$ は $R_\p$ のただ一つの極大イデアルである．$R_\p$ の素イデアルは，$\p$ に含まれる $R$ の素イデアル $\q$ と $\q\mapsto\q R_\p$ により一対一に対応する．
\end{myprop}

\begin{proof}
$x=a/s\notin\p R_\p$ なら $a\notin\p$（もし $a/s=b/t$，$b\in\p$ なら $u(at-bs)=0$，$uat\in\p$，$u,t\notin\p$ より $a\in\p$），
よって $s/a\in R_\p$ が逆元であり $x$ は単元．つまり $\p R_\p$ の外はすべて単元なので，任意の真のイデアルは $\p R_\p$ に含まれる．
$\p R_\p\ne R_\p$ は命題 \ref{prop:ch06-loc-prime} から分かる．素イデアルの対応は同命題で $S=R\setminus\p$ とおけばよい（$\q\cap S=\emptyset\iff\q\subset\p$）．
\end{proof}

$R_\p$ の剰余体 $\kappa(\p)=R_\p/\p R_\p$ は $\Frac(R/\p)$ と同型である．
例えば $R=\Z$，$\p=(p)$ なら $\kappa(\p)=\F_p$，$\p=(0)$ なら $\kappa(\p)=\Q$．

\begin{figure}[H]
\centering
\begin{tikzpicture}[scale=0.9,every node/.style={font=\small}]
  \node (R) at (0,0) {$R$ の素イデアル};
  \node (Rp) at (7,0) {$R_\p$ の素イデアル};
  \draw[rounded corners,thick] (-2.2,-3.2) rectangle (2.2,-0.5);
  \draw[rounded corners,thick,fill=blue!8] (-1.9,-3.0) rectangle (1.2,-0.9);
  \node at (-0.4,-1.2) {$\q\subset\p$ の部分};
  \node at (-0.4,-2.0) {$\p$};
  \node at (1.7,-2.0) {$\cdots$};
  \draw[rounded corners,thick,fill=blue!8] (5.4,-3.0) rectangle (8.6,-0.9);
  \node at (7,-2.0) {$\p R_\p$（唯一の極大）};
  \draw[->,thick] (1.3,-2.3)--(5.3,-2.3) node[midway,above]{$\q\mapsto\q R_\p$};
\end{tikzpicture}
\caption{$R_\p$ の素イデアルは $\p$ に含まれる素イデアルだけが生き残る．}
\end{figure}

\subsection{極大イデアルによる局所化}

$\m$ が極大イデアルの場合が，幾何学的に最も重要な「点における局所化」である．
$X$ をアフィン多様体，$p\in X$，$\m_p=\{f\in k[X]\mid f(p)=0\}$ とすると
\[
  k[X]_{\m_p}=\{\,g/h\mid g,h\in k[X],\ h(p)\ne0\,\}\subset k(X).
\]
これは\ref{subsec:ch06-germ}小節で関数の芽として定義した $\OO_{X,p}$ と一致する（\ref{subsec:ch06-OXp}小節で改めて述べる）．
$R_\m$ の剰余体は $R/\m$ であり，点の場合は $k[X]/\m_p\cong k$ である（芽 $\mapsto p$ での値）．

\begin{myex}{整数の局所化 $\Z_{(p)}$}{ch06-Zp}
素数 $p$ に対し
\[
  \Z_{(p)}=\Bigl\{\frac ab\ \Big|\ a,b\in\Z,\ p\nmid b\Bigr\}\subset\Q .
\]
単元は $p\nmid a$ なる $a/b$，唯一の極大イデアルは $p\Z_{(p)}$．$\Z_{(p)}$ のイデアルは $(p^n)$（$n\ge0$）と $0$ のみであり，
すべての非零元は $u p^n$（$u$ 単元）と一意に書ける．$n=v_p(x)$ は $p$ 進付値である．
「素数 $p$ の近くでの整数」を見ているといえ，算術幾何では $\Spec\Z$ の点 $(p)$ における局所環である．
\end{myex}

\subsection{局所化と点の周辺の代数}

$\OO_{X,p}=k[X]_{\m_p}$ の素イデアルは，命題 \ref{prop:ch06-Rp-local} により $\m_p$ に含まれる $k[X]$ の素イデアル，
すなわち\emph{点 $p$ を通る既約閉部分多様体}と一対一に対応する．
$p$ を通らない部分多様体の情報は局所化によって消える．これが「点の周辺だけを見る」ことの代数的な実現である．

\begin{myex}{$k[x,y]$ の原点での局所化}{ch06-A2-origin}
$R=k[x,y]$，$\m=(x,y)$ とする．$R_\m$ の素イデアルは，原点を通る既約閉部分集合に対応する：
\begin{itemize}
  \item $0$ $\leftrightarrow$ $\A^2$ 全体，
  \item $(f)$，$f$ 既約で $f(0,0)=0$ $\leftrightarrow$ 原点を通る既約曲線，
  \item $\m R_\m$ $\leftrightarrow$ 原点．
\end{itemize}
直線 $x=1$ に対応する素イデアル $(x-1)$ は $x-1\notin\m$ なので $R_\m$ では単元を含み，消える．
\end{myex}

\begin{myex}{node の局所環と既約成分の比較}{ch06-node-local}
$X=\Zero(xy)$，$R=k[X]=k[x,y]/(xy)$ を考える．
\begin{itemize}
  \item $p=(1,0)$ のとき $\m_p=(x-1,y)$．$x\notin\m_p$ なので $x$ は $R_{\m_p}$ で単元となり，$xy=0$ から $y=0$．
        よって $R_{\m_p}\cong k[x]_{(x-1)}$：$p$ の近くでは $X$ は $x$ 軸だけに見える．
  \item $p=(0,0)$ のとき $x,y$ はどちらも単元でなく，$R_{\m_p}$ には零因子 $x,y$ が残る．$R_{\m_p}$ の極小素イデアルは $(x)$，$(y)$ の二つで，
        原点では二つの成分が両方見えている．
\end{itemize}
局所化は「その点を通る成分だけを残す」操作であることがよく分かる．
\end{myex}

\exprob{prob:ch06-3-1}
$R=k[x,y]/(y^2-x^2(x+1))$（結節曲線）とし，$p=(0,0)$，$q=(-1,0)$ とする．$R_{\m_q}$ が離散付値環（極大イデアルが $0$ でない単項イデアルである Noether 局所整域）であることを示せ．
また $R_{\m_p}$ の極大イデアルは単項でないことを示せ（ヒント：$\m/\m^2$ の次元を計算せよ）．

% ----------------------------------------------------------
\section{局所環}
% ----------------------------------------------------------

\subsection{局所環の定義}

\begin{mydef}{局所環}{ch06-local-ring}
環 $R$ がただ一つの極大イデアル $\m$ をもつとき，$R$ を\textbf{局所環}といい，$(R,\m)$ あるいは $(R,\m,\kappa)$ と書く．
ここで $\kappa=R/\m$ を\textbf{剰余体}という．
局所環の間の環準同型 $\phi:(R,\m)\to(R',\m')$ が $\phi(\m)\subset\m'$ を満たすとき，\textbf{局所準同型}という．
\end{mydef}

前節の結果から $R_\p$ は常に局所環である．「局所化」という名前は，局所化によって局所環が得られることに由来する．

局所準同型の条件は自然である．射 $\varphi:X\to Y$，$q=\varphi(p)$ に対し，引き戻し $\varphi^*:\OO_{Y,q}\to\OO_{X,p}$ は
$q$ で消える関数を $p$ で消える関数に写すので局所準同型である．

\subsection{極大イデアルの一意性}

局所環であることは次のように特徴付けられる．

\begin{myprop}{局所環の特徴付け}{ch06-local-char}
環 $R$ について次は同値である．
\begin{enumerate}
  \item $R$ は局所環である．
  \item $R$ の非単元全体 $\m=R\setminus R^\times$ はイデアルである．
  \item ある真のイデアル $\m$ が存在して，$\m$ に属さない元はすべて単元である．
  \item ある極大イデアル $\m$ が存在して，任意の $a\in\m$ に対し $1+a$ は単元である．
\end{enumerate}
\end{myprop}

\begin{proof}
(1)$\Rightarrow$(2)：任意の非単元 $x$ は $(x)\neq R$ なのである極大イデアルに含まれ，それは $\m$ である．よって $R\setminus R^\times\subset\m$．
$\m$ は単元を含まないので等号．
(2)$\Rightarrow$(3)：明らか．
(3)$\Rightarrow$(1)：任意の真のイデアル $I$ は単元を含まないので $I\subset\m$．よって $\m$ は唯一の極大イデアル．
(1)$\Rightarrow$(4)：$1+a$ が単元でなければある極大イデアル，すなわち $\m$ に属し，$1=(1+a)-a\in\m$ となり矛盾．
(4)$\Rightarrow$(3)：$x\notin\m$ とすると，$\m$ の極大性から $\m+(x)=R$，$1=a+rx$（$a\in\m$）．$rx=1-a$ は (4) より単元なので $x$ も単元．
\end{proof}

(3) は局所環であることを確かめる実用的な判定法である．例えば $R_\p$ では「$\p R_\p$ の外はすべて単元」であった（命題 \ref{prop:ch06-Rp-local}）．
幾何学的には「点 $p$ で値が $0$ でない関数は，$p$ の近くで逆数をとれる」ということであり，
局所環とは\emph{点で消えない関数がすべて可逆になっている}関数環である．

\subsection{局所環の基本性質}

局所環の基本的な道具が中山の補題である．これは第7章で Zariski 接空間 $\m/\m^2$ を扱う際に本質的に用いる．

\begin{mylem}{中山の補題}{ch06-nakayama}
$(R,\m)$ を局所環，$M$ を有限生成 $R$ 加群とする．
\begin{enumerate}
  \item $\m M=M$ ならば $M=0$．
  \item $N\subset M$ を部分加群とし，$M=N+\m M$ ならば $M=N$．
  \item $m_1,\dots,m_r\in M$ の像が $\kappa$ ベクトル空間 $M/\m M$ を張るならば，$m_1,\dots,m_r$ は $M$ を生成する．
\end{enumerate}
\end{mylem}

\begin{proof}
(1) $M\neq 0$ として，生成元の個数が最小な生成系 $m_1,\dots,m_r$ をとる．$m_r\in\m M$ より $m_r=\sum_i a_im_i$（$a_i\in\m$），
$(1-a_r)m_r=\sum_{i<r}a_im_i$．命題 \ref{prop:ch06-local-char} (4) より $1-a_r$ は単元なので $m_r$ は $m_1,\dots,m_{r-1}$ で書け，最小性に矛盾．
(2) $M/N$ に (1) を適用する．(3) $N=\sum Rm_i$ に (2) を適用する．
\end{proof}

\begin{mycor}{極小生成系}{ch06-min-gen}
$(R,\m,\kappa)$ を Noether 局所環とすると，$\m$ の極小な生成系の元の個数は $\dim_\kappa\m/\m^2$ に等しい．
\end{mycor}

この数 $\dim_\kappa\m/\m^2$ は（$R=\OO_{X,p}$ のとき）点 $p$ における Zariski 接空間の次元であり，第7章の主役となる．

その他の基本性質として，
\begin{itemize}
  \item 局所環の剰余体 $\kappa$ は体であり，局所環 $R$ から体への準同型で核が $\m$ を含むものは $\kappa$ を経由する．
  \item 局所環の冪等元は $0,1$ のみ（$e(1-e)=0$ で，$e,1-e$ の一方は単元）．幾何学的には「局所的には連結」である．
  \item 局所環 $R$ の商 $R/I$（$I\neq R$）も局所環であり，局所化 $R_\p$ も局所環である．
\end{itemize}

\subsection{局所環の例}

\begin{myex}{局所環の例}{ch06-local-ex}
\begin{enumerate}
  \item 体 $k$：極大イデアルは $0$．
  \item $\Z_{(p)}$（例 \ref{ex:ch06-Zp}）：剰余体 $\F_p$．
  \item $k[x]_{(x)}=\{f/g\mid g(0)\neq0\}$：原点で分母が消えない有理関数．極大イデアルは $(x)$，剰余体 $k$（$f/g\mapsto f(0)/g(0)$）．
        非零元は $x^n u$（$u$ 単元）と一意に書け，$n$ は原点での零点の位数である．
  \item 形式的冪級数環 $k[[x_1,\dots,x_n]]$：$f=\sum a_\alpha x^\alpha$ が単元であるのは定数項 $a_0\neq0$ のとき，そしてそのときに限る．
        実際 $f=a_0(1-g)$（$g$ は定数項 $0$）なら $f^{-1}=a_0^{-1}\sum_{m\ge0}g^m$ は形式的に収束する．
        よって非単元全体 $(x_1,\dots,x_n)$ はイデアルであり，命題 \ref{prop:ch06-local-char} (2) により局所環．
  \item 収束冪級数環 $\C\{z_1,\dots,z_n\}$（原点の近傍で収束する冪級数の全体，すなわち正則関数の芽）．定数項 $\ne0$ なら逆数も原点近傍で正則．
  \item $k[x]/(x^n)$：唯一の素イデアル $(x)$ をもつ局所環（Artin 局所環）．
  \item 非例：$k[x]$ は局所環でない（$(x)$ と $(x-1)$ が異なる極大イデアル）．$\Z$ も同様．
\end{enumerate}
\end{myex}

(3)(4)(5) の間には包含
\[
  k[x]\subset k[x]_{(x)}\subset \C\{x\}\subset\C[[x]]\qquad(k=\C)
\]
がある．$1/(1-x)\in k[x]_{(x)}$ は $\sum x^n$ と展開され，局所化の元は「原点での Taylor 展開」によって冪級数とみなせる．
代数的な局所環 $k[x]_{(x)}$ の元は有理関数であり，$\C\{x\}$ の元（例えば $e^x$）の大部分は含まない．
$\C[[x]]$ は $k[x]_{(x)}$ の $\m$ 進完備化 $\varprojlim k[x]_{(x)}/(x^n)$ である．

\subsection{多様体の点に付随する局所環}
\label{subsec:ch06-OXp}

\begin{mydef}{点における局所環}{ch06-OXp}
アフィン多様体 $X$ と点 $p\in X$ に対し，
\[
  \OO_{X,p}:=k[X]_{\m_p}=\{\,g/h\mid g,h\in k[X],\ h(p)\neq0\,\}
\]
を $X$ の $p$ における\textbf{局所環}という．その極大イデアルを $\m_{X,p}=\m_p\OO_{X,p}$ と書く．
\end{mydef}

\begin{myprop}{芽としての局所環}{ch06-OXp-germ}
$\OO_{X,p}$ は，\ref{subsec:ch06-germ}小節で定義した $p$ における正則関数の芽のなす環と自然に同型である．
$\OO_{X,p}$ は剰余体 $k$ の局所環で，剰余写像は芽の $p$ での値 $\varphi\mapsto\varphi(p)$ である．
\end{myprop}

\begin{proof}
$g/h$（$h(p)\neq0$）に，$D(h)$ 上の正則関数 $g/h$ の芽を対応させる．これは well-defined で（同値な分数は $k(X)$ で等しく，共通の開集合上で一致する），
単射である（芽が $0$ なら空でない開集合上で $g=0$，一致の定理より $g=0$）．
全射性：芽 $[U,f]$ をとると，正則関数の定義から $p$ の近傍 $V$ で $f=g/h$，$h$ は $V$ 上で $0$ でない．とくに $h(p)\neq0$．
局所環であることは命題 \ref{prop:ch06-Rp-local}．
\end{proof}

\begin{myex}{いくつかの $\OO_{X,p}$}{ch06-OXp-ex}
\begin{enumerate}
  \item $\OO_{\A^1,0}=k[x]_{(x)}$．
  \item $\OO_{\A^n,0}=k[x_1,\dots,x_n]_{(x_1,\dots,x_n)}$．
  \item 放物線 $X=\Zero(y-x^2)$ の原点：$\OO_{X,0}\cong k[t]_{(t)}$（$X\cong\A^1$ なので）．
  \item 尖点曲線 $C=\Zero(y^2-x^3)$ の原点：$\OO_{C,0}\cong k[t^2,t^3]_{(t^2,t^3)}$．これは $k[t]_{(t)}$ と同型でない（$\m/\m^2$ が2次元，第7章）．
        原点以外の点 $p$ では $\OO_{C,p}\cong\OO_{\A^1,t}$．実際 $t=y/x$ は $x\neq0$ の近傍で正則なので，$\varphi:\A^1\to C$ は $C\setminus\{0\}$ 上で同型である．
\end{enumerate}
\end{myex}

(4) は，第5章で見た「尖点曲線は $\A^1$ と同相だが同型でない」ことの局所的な説明になっている：
同型でないことの原因は原点の局所環だけにある．

\exprob{prob:ch06-4-1}
$k[[x]]$ のイデアルは $0$ と $(x^n)$（$n\ge0$）に限ることを示せ．また $k[x]_{(x)}$ についても同じ主張が成り立つことを示せ．

\exprob{prob:ch06-4-2}
$(R,\m)$ を局所環とする．$R$ の元 $e$ が $e^2=e$ を満たせば $e=0$ または $e=1$ であることを示せ．
これを用いて，$\Zero(xy)\setminus\{0\}$ のような非連結な空間の座標環は局所環ではありえないことを説明せよ．

% ----------------------------------------------------------
\section{局所環と幾何学}
% ----------------------------------------------------------

\subsection{点の近傍における関数}

$\OO_{X,p}$ は「$p$ の近傍で定義された関数を，近傍を十分小さくしても区別できるものだけ区別したもの」であった．
幾何学的にいくつかの基本的な事実を確認する．

\begin{itemize}
  \item \textbf{値と極大イデアル．}芽 $\varphi\in\OO_{X,p}$ について，$\varphi\in\m_{X,p}\iff\varphi(p)=0$．
        $\varphi(p)\ne0$ なら $\varphi$ は $\OO_{X,p}$ で可逆（$p$ の近傍で $1/\varphi$ がとれる）．
  \item \textbf{消える位数．}$\varphi\in\m_{X,p}^n\setminus\m_{X,p}^{n+1}$ のとき，$\varphi$ は $p$ で「ちょうど $n$ 次で消える」と考えられる．
        $X=\A^n$ なら，これは Taylor 展開の最低次の項が $n$ 次であることを意味する．
  \item \textbf{大域関数の復元．}既約な $X$ で $k[X]=\bigcap_{p\in X}\OO_{X,p}$（$k(X)$ の中で）．
        これは第5章の定理（$\OO_X(X)=k[X]$）の言い換えである：すべての点で正則な有理関数は多項式関数である．
\end{itemize}

最後の性質は，「局所的な情報をすべて集めると大域的な情報が復元できる」ことの最初の例であり，
第11章の層の考え方の原型である．代数的に言えば，整域 $R$ について $R=\bigcap_{\m}R_\m$（$\m$ は極大イデアル全体）が成り立つ：
$x\in\bigcap R_\m$ に対し $J=\{a\mid ax\in R\}$ がどの極大イデアルにも含まれないので $J=R$．

\subsection{局所的な代数構造}

局所環 $(\OO_{X,p},\m)$ から次のような代数的構造が取り出せる．
\begin{enumerate}
  \item \textbf{剰余体} $\OO_{X,p}/\m\cong k$：関数の値．
  \item \textbf{余接空間} $\m/\m^2$：関数の「1次の部分」．$X=\A^n$，$p=0$ なら $\m/\m^2$ は $\bar x_1,\dots,\bar x_n$ を基底とする $n$ 次元空間で，
        $f\mapsto\sum_i\frac{\partial f}{\partial x_i}(0)\,\bar x_i$ は微分 $df_0$ にほかならない．その双対 $(\m/\m^2)^*$ が Zariski 接空間（第7章）．
  \item \textbf{随伴次数環} $\mathrm{gr}_\m=\bigoplus_{n\ge0}\m^n/\m^{n+1}$：Taylor 展開の各次の部分．幾何学的には接錐に対応する．
  \item \textbf{完備化} $\widehat{\OO}_{X,p}=\varprojlim_n\OO_{X,p}/\m^n$：形式的 Taylor 展開全体．非特異点 $p$（次元 $d$）では $\widehat{\OO}_{X,p}\cong k[[t_1,\dots,t_d]]$ となる．
\end{enumerate}
$\OO_{X,p}/\m^{n+1}$ は「$p$ での $n$ 次までの Taylor 展開（ジェット）」の環である．
有限次元の $k$ 代数であり，$p$ の「$n$ 次の無限小近傍」の関数環とみなせる（第12章でスキームとして解釈する）．

\begin{myex}{尖点曲線の局所構造}{ch06-cusp-local}
$C=\Zero(y^2-x^3)$，$p=0$，$\m=(\bar x,\bar y)$ とする．$\m^2=(\bar x^2,\bar x\bar y,\bar y^2)$ で，$y^2-x^3$ は2次以上の項しかないので，
$\m/\m^2$ は $\bar x,\bar y$ を基底とする2次元空間となる．曲線（1次元）なのに余接空間が2次元であることが，原点が特異点であることを表す．
さらに $\mathrm{gr}_\m\cong k[x,y]/(y^2)$（定義式の最低次の部分）であり，接錐は二重直線 $y^2=0$ である．
\end{myex}

\subsection{局所的な幾何の解析}

局所環は「点の近くで $X$ がどう見えるか」を調べる道具である．
二つの点 $p\in X$，$q\in Y$ について，$\OO_{X,p}\cong\OO_{Y,q}$ であれば，$p$ と $q$ のある開近傍が同型であることが知られている
（局所環は Zariski 局所的な同型類を決める）．
さらに粗く，完備化 $\widehat\OO_{X,p}\cong\widehat\OO_{Y,q}$ のとき「$p$ と $q$ は解析的に同型」という．
例えば $k=\C$ では，これは Euclid 位相での小さな近傍が複素解析的に同型であることとほぼ同じである．

\begin{myex}{結節点の枝：局所環と完備化の違い}{ch06-node-branch}
結節曲線 $N=\Zero(y^2-x^2(x+1))$ の原点を考える．
$y^2-x^2(x+1)$ は $k[x,y]$ で既約なので，$\OO_{N,0}$ は整域であり，Zariski 局所的には $N$ は既約で「1本」に見える．
ところが完備化すると，$\sqrt{1+x}=1+\tfrac12x-\tfrac18x^2+\cdots\in k[[x]]$ を用いて
\[
  y^2-x^2(1+x)=\bigl(y-x\sqrt{1+x}\bigr)\bigl(y+x\sqrt{1+x}\bigr)
\]
と分解し，$\widehat\OO_{N,0}\cong k[[x,y]]/(uv)$（$u,v=y\mp x\sqrt{1+x}$）は整域でない．
すなわち\emph{形式的（解析的）には原点で2本の枝が交わっている}．これは $\Zero(xy)$ の原点（例 \ref{ex:ch06-node-local}）と解析的に同型である．
Zariski 近傍は大きすぎて2本の枝を分離できない（遠くで繋がっている）が，無限小の近傍では分離できる．
\end{myex}

\begin{figure}[H]
\centering
\begin{tikzpicture}[scale=1.1]
  \draw[thick,teal,domain=-1.35:1.35,samples=80] plot ({\x*\x-1},{\x*(\x*\x-1)});
  \draw[dashed,gray] (0,0) circle (0.45);
  \fill (0,0) circle (1.5pt);
  \node[teal] at (0.9,0.9) {$N$};
  \draw[->] (0.5,0.1)--(2.2,0.6);
  \begin{scope}[xshift=3.3cm,yshift=0.6cm]
    \draw[dashed,gray] (0,0) circle (0.8);
    \draw[thick,teal] (-0.56,-0.56)--(0.56,0.56);
    \draw[thick,teal] (-0.56,0.56)--(0.56,-0.56);
    \fill (0,0) circle (1.5pt);
    \node[font=\small,align=center] at (0,-1.3) {無限小近傍：\\ 2本の枝 $uv=0$};
  \end{scope}
\end{tikzpicture}
\caption{結節点の近傍．Zariski 局所的には既約だが，完備化すると2本の枝に分かれる．}
\label{fig:ch06-node}
\end{figure}

\subsection{局所環による点の構造の記述}

点 $p$ の「種類」は局所環の代数的性質で記述される．本ノートで扱うものをまとめておく．

\begin{center}
\renewcommand{\arraystretch}{1.35}
\begin{tabular}{l|l}
点 $p\in X$ の幾何 & 局所環 $\OO_{X,p}$ の代数\\ \hline
$p$ を通る既約成分 & $\OO_{X,p}$ の極小素イデアル\\
$p$ を通る部分多様体 & $\OO_{X,p}$ の素イデアル\\
$p$ での $X$ の次元 & $\OO_{X,p}$ の Krull 次元（第7章）\\
接空間 $T_pX$ & $(\m/\m^2)^*$（第7章）\\
$p$ は非特異点 & $\dim_k\m/\m^2=\dim\OO_{X,p}$（正則局所環，第7章）\\
接錐 & 随伴次数環 $\mathrm{gr}_\m\OO_{X,p}$\\
解析的な枝 & $\widehat\OO_{X,p}$ の極小素イデアル\\
曲線上の点での零点・極の位数 & 離散付値（$\OO_{X,p}$ が DVR のとき）
\end{tabular}
\end{center}

とくに曲線の非特異点では $\OO_{X,p}$ は離散付値環（DVR）となり，局所パラメータ $t$（$\m=(t)$）を用いて
任意の $f\in k(X)^\times$ を $f=t^nu$（$u$ は単元）と一意に書ける．$n=\mathrm{ord}_p(f)$ は $f$ の $p$ での零点（$n>0$）・極（$n<0$）の位数であり，
複素解析における Laurent 展開の最低次数と同じものである．これは因子や Riemann--Roch の定理（後の章）の基礎となる．

\begin{physbox}[局所パラメータと Laurent 展開]
複素平面上の有理型関数 $f(z)$ を $z=a$ のまわりで Laurent 展開したときの最低次数 $n$ は，留数計算や共形場理論の
演算子積展開（OPE）$A(z)B(w)\sim\sum_n\frac{C_n(w)}{(z-w)^n}$ で現れる特異性の次数である．
代数幾何ではこれが局所環 $\OO_{X,p}$（DVR）の付値 $\mathrm{ord}_p$ として定式化される．
局所パラメータ $t$ の選び方は単元倍の不定性をもつが，$\mathrm{ord}_p$ はそれによらない．
これは座標変換で極の次数が変わらないことに対応している．
\end{physbox}

\exprob{prob:ch06-5-1}
$X=\A^2$，$p=0$ とし，$f=x^2+y^3$，$g=x+y^2$ とする．$f,g$ がそれぞれ $\m^n\setminus\m^{n+1}$ に属する $n$ を求めよ．
また $\OO_{X,p}/(f)$ と $\OO_{X,p}/(g)$ の $\m/\m^2$ の次元を計算し，$\Zero(f)$，$\Zero(g)$ のどちらが原点で特異かを判定せよ．

% ----------------------------------------------------------
\section{局所化の物理的背景}
\label{sec:ch06-phys}
% ----------------------------------------------------------

\subsection{局所的な記述と大域的な記述}

物理学でも数学でも，「局所的に記述して大域的に貼り合わせる」という戦略は共通している．
多様体上のベクトル束はチャートごとに自明化し，遷移関数で貼り合わせる．ゲージ場はパッチごとにゲージポテンシャル $A_\mu$ を与え，
重なりでゲージ変換によって繋ぐ（Wu--Yang の磁気単極子）．
代数幾何では，アフィン開集合 $D(f)$ への制限が局所化 $k[X]\to k[X]_f$ であり，
点への「無限に小さな」制限が $k[X]\to\OO_{X,p}$ である．

ただし代数幾何の局所化には二つの段階があることを強調しておく：
\begin{center}
\renewcommand{\arraystretch}{1.3}
\begin{tabular}{l|l|l}
操作 & 代数 & 幾何的意味\\ \hline
開集合への制限 & $R\to R_f$ & 有限個の超曲面を除いた，大きな開集合\\
点での芽 & $R\to R_\m$ & すべての開近傍の極限（帰納極限）\\
無限小近傍 & $R\to\widehat{R_\m}$ & 形式的 Taylor 展開
\end{tabular}
\end{center}
$\OO_{X,p}=\varinjlim_{p\in D(f)}k[X]_f$ と書けることに注意すると，「点での局所化」は「開集合への制限の極限」として理解できる．

\subsection{場の理論における局所性との区別}

物理の「局所性（locality）」は主に次の意味で使われる．
\begin{enumerate}
  \item \textbf{作用の局所性：} 作用が $S=\int d^dx\,\mathcal L(\phi(x),\partial\phi(x),\dots)$ と，各点での場とその\emph{有限階}の微分の関数の積分で書けること．
  \item \textbf{微視的因果律（局所可換性）：} 空間的に離れた点 $x,y$ に対して $[\mathcal O(x),\mathcal O(y)]=0$．
        代数的場の量子論では，時空の領域 $\mathcal R$ ごとに局所観測量の代数 $\mathfrak A(\mathcal R)$ を対応させ，
        空間的に離れた領域の代数が可換であることを公理とする（Haag--Kastler）．
  \item \textbf{クラスター分解：} 遠く離れた実験の結果は互いに独立になる．
\end{enumerate}
これらと代数幾何の「局所化・局所環」の関係を整理しよう．

\paragraph{(a) 「局所」の空間的スケールが全く違う．}
物理の局所性は Euclid 的（あるいは因果的）な距離に基づき，「任意に小さな領域」を考えることが前提である．
これに対し Zariski 開集合は有限個の超曲面を除いた巨大な集合であり，距離の概念は存在しない．
代数幾何の「$p$ の近く」とは「$p$ を除外しない程度に小さい」という意味でしかなく，
既約な $X$ では任意の二つの開近傍は交わる．したがって「空間的に離れた二つの領域」という概念がそもそも Zariski 位相では表現できない．

\paragraph{(b) 関数の剛性．}
物理の局所的な場は，コンパクトな台をもつ試験関数 $f$ で「ぼかした」演算子 $\phi(f)=\int f(x)\phi(x)\,d^dx$ として扱われ，
台が互いに重ならない試験関数を自由にとれることが局所性の定式化に不可欠である（bump function の存在）．
代数幾何の正則関数・有理関数は一致の定理により，芽が決まれば全体が決まる（\ref{subsec:ch06-germ}小節）．
$\OO_{X,p}$ の元は「$p$ の近くの情報」でありながら，実は定義域全体での値を決めてしまう．
コンパクト台をもつ多項式関数は $0$ しかない．
この意味で，代数幾何の局所環は物理の局所観測量の代数とはまったく性質が異なり，むしろ\emph{解析接続}の世界に近い．

\paragraph{(c) 因果構造の不在．}
微視的因果律は Lorentz 計量の光円錐構造に依存するが，代数多様体の局所環にはそのような構造はない．
さらに局所環は可換環であり，量子論の非可換な局所代数とは代数的な型も異なる．
「局所化」という言葉が共通しているのは偶然ではない（どちらも「ある場所の近くで意味をもつ量だけに注目する」）が，
\emph{代数幾何の局所化は因果性や相対論的局所性とは独立な，純粋に環論的な操作である}と理解すべきである．

\paragraph{(d) 共通点：「局所から大域へ」．}
一方で，「局所的なデータの整合的な貼り合わせとして大域的な対象を記述する」という構造は共通しており，
代数的場の量子論の「領域 $\mathcal R\mapsto\mathfrak A(\mathcal R)$」という対応（ネット）は，
開集合 $U\mapsto\OO_X(U)$ という対応（層，第11章）と形式的に似ている．
両者の違いは，前者が非可換代数で因果的な公理を課すのに対し，後者は可換環で貼り合わせの公理（層の公理）を課す点にある．

\begin{physbox}[有効場の理論と完備化]
有効場の理論（EFT）では，真空 $\phi=\phi_0$ のまわりでポテンシャルを $V(\phi)=\sum_n\frac{1}{n!}V^{(n)}(\phi_0)(\phi-\phi_0)^n$ と展開し，
低エネルギーでは低次の項だけを残す．これは代数的には局所環 $\OO_{X,\phi_0}$ の $\m$ 進フィルトレーション
$\OO\supset\m\supset\m^2\supset\cdots$ に沿って「$n$ 次まで」で打ち切ること，すなわち $\OO_{X,p}/\m^{n+1}$（ジェット）を見ることに対応する．
真空の安定性や質量スペクトル（2次の項 $\m^2/\m^3$），相互作用（3次以上）といった情報は局所環の随伴次数環から読み取れる．
摂動展開は形式的冪級数 $\widehat\OO$ の中の計算であり，その「収束」は代数幾何の局所環では問われない別の問題である
（漸近級数・非摂動効果）．このように，物理の「局所」には\emph{空間的局所性}と\emph{場の配位空間（あるいはパラメータ空間）上の点のまわりの展開}という
二つの異なる意味があり，代数幾何の局所環が直接対応するのは主に後者である．
\end{physbox}

\subsection{局所的な情報を扱う数学的手法}

本章で導入した手法を，その使いどころとともにまとめておく．
\begin{itemize}
  \item \textbf{局所化 $S^{-1}R$：} 分母を許す普遍的な操作．開集合への制限（$R_f$），点での芽（$R_\p$），関数体（$\Frac R$）を統一的に扱う．
        素イデアルの対応（命題 \ref{prop:ch06-loc-prime}）により「どの部分多様体が見えるか」が管理できる．
  \item \textbf{局所環と中山の補題：} 点での接空間（第7章），重複度，交点数（第9章）など点の不変量の計算の基礎．
  \item \textbf{完備化：} 形式的・解析的な局所構造．特異点の枝の分離，Taylor 展開による計算．
  \item \textbf{層（第11章）：} 開集合ごとの局所化 $U\mapsto\OO_X(U)$ を整合的に束ねたもの．局所環はその茎 $\OO_{X,p}=\varinjlim_{U\ni p}\OO_X(U)$ である．
  \item \textbf{スキーム（第12章）：} 素イデアルでの局所化 $R_\p$ をすべての素イデアルについて考え，$\Spec R$ を局所環つき空間として構成する．
\end{itemize}

物理への応用の観点からは，局所環は\emph{特異点の近くの物理}を調べる道具として重要である．
弦理論のコンパクト化で特異点（オービフォールド点，コニフォールド点など）が現れるとき，
その近傍の物理（新たな無質量状態，ゲージ対称性の増大など）は特異点の局所環・完備化によって分類される局所構造でほぼ決まる．
例えばコニフォールド特異点 $xy-zw=0$ の局所環の構造は，その点の近くでの D ブレーンの物理を決める．
これらは第7章の特異点論，および第IV部以降で再び取り上げる．

\exprob{prob:ch06-6-1}
$\R$ 上の $C^\infty$ 関数の原点における芽のなす環 $C^\infty_0$ は局所環であることを示せ．
さらに，$e^{-1/x^2}$（$x\neq0$），$0$（$x=0$）で定まる関数の芽は $0$ でないが，その Taylor 展開（$\m^\infty=\bigcap_n\m^n$ を法とする像）は $0$ であることを確かめ，
代数的な局所環 $k[x]_{(x)}$ では $\bigcap_n\m^n=0$ となる（Krull の交叉定理）こととの違いを論ぜよ．

\section*{問題の解答}
\addcontentsline{toc}{section}{問題の解答}

\exsol{prob:ch06-1-1}
\textbf{芽 $1/(1-x)$ と $\sum x^n$ の関係．}
$\OO_{\A^1,0}=k[x]_{(x)}=\{f/g\mid g(0)\ne0\}$ であり，芽 $[\A^1\setminus\{1\},1/(1-x)]$ はこの環の元 $1/(1-x)$ に対応する（分母 $1-x$ は原点で $1$）．
一方，形式的冪級数環 $k[[x]]$ では定数項が $0$ でない元は可逆である（例\ref{ex:ch06-local-ex}(4)）．
したがって局所化の普遍性（定理\ref{thm:ch06-universal}）により，包含 $k[x]\subset k[[x]]$ は環準同型
\[
  T\colon k[x]_{(x)}\longrightarrow k[[x]],\qquad f/g\longmapsto f\cdot g^{-1}
\]
に一意に延びる．これは原点での Taylor 展開を与える写像であり，$k[[x]]$ が整域なので単射である（$fg^{-1}=0\Rightarrow f=0$）．
$(1-x)\sum_{n\ge0}x^n=1$ が $k[[x]]$ で成り立つので
\[
  T\Bigl(\frac1{1-x}\Bigr)=\sum_{n\ge0}x^n .
\]
すなわち $\sum x^n$ は芽 $1/(1-x)$ の原点での形式的 Taylor 展開である．さらに $\m$ を法とする各段階で見ると
\[
  \frac1{1-x}\equiv1+x+\cdots+x^n\pmod{\m^{n+1}}
\]
（$1/(1-x)-(1+\cdots+x^n)=x^{n+1}/(1-x)\in\m^{n+1}$）であり，冪級数は「全ての有限次ジェットを集めたもの」である．

ただし両者の性格は異なる．$k=\C$ のとき級数 $\sum x^n$ は $|x|<1$ でしか収束しないが，芽 $1/(1-x)$ は Zariski 開集合 $\A^1\setminus\{1\}$ 全体の正則関数を定めている．
代数的な芽は（一致の定理により）定義域全体での関数を決めてしまうのに対し，冪級数は原点の無限小近傍の情報だけを見ている．
$T$ は単射なので情報は失われないが，$T$ は全射ではない（例えば $e^x=\sum x^n/n!$ は有理関数の展開ではない）．

\textbf{多項式関数にならない元の例．}
$1/(1-x)\in\OO_{\A^1,0}$ 自身がその例である．もし $k[x]$ の元 $p(x)$ と（$k(x)$ の元として）一致すれば $(1-x)p(x)=1$ となるが，左辺は $p\ne0$ なら1次以上なので矛盾．
幾何学的には，$1/(1-x)$ は $x=1$ で極をもつので $\A^1$ 全体には延びない．
同様に $1/(x-a)$（$a\ne0$）や $1/(x^2+1)$（$k=\C$）なども原点では正則だが多項式関数ではない．

\exsol{prob:ch06-2-1}
$R=k[x,y]/(xy)$，$f=x+y$ とし，$R_f=S^{-1}R$ を計算する．答えは
\[
  R_{x+y}\cong k[x,x^{-1}]\times k[y,y^{-1}]
\]
である．

\textbf{冪等元．}
$R$ では $xy=0$ なので $x(x+y)=x^2$，$y(x+y)=y^2$．$R_{x+y}$ において
\[
  e=\frac{x}{x+y},\qquad 1-e=\frac{y}{x+y}
\]
とおくと，$e^2=\frac{x^2}{(x+y)^2}=\frac{x(x+y)}{(x+y)^2}=e$，$e(1-e)=\frac{xy}{(x+y)^2}=0$ である．$e$ は「$x$ 軸上で $1$，$y$ 軸上で $0$」の関数である．

\textbf{同型の構成．}
環準同型 $\phi\colon R\to k[x,x^{-1}]\times k[y,y^{-1}]$ を $\phi(h)=(h(x,0),\ h(0,y))$ で定める（$xy\mapsto(0,0)$ なので well-defined）．
$\phi(x+y)=(x,y)$ は右辺の単元なので，普遍性（定理\ref{thm:ch06-universal}）により $\tilde\phi\colon R_{x+y}\to k[x^{\pm1}]\times k[y^{\pm1}]$ が得られる．
\begin{itemize}
  \item \emph{全射．}$\tilde\phi(e)=(1,0)$，$\tilde\phi(1-e)=(0,1)$，$\tilde\phi(x)=(x,0)$，$\tilde\phi(y)=(0,y)$，$\tilde\phi(1/(x+y))=(x^{-1},y^{-1})$ であり，
        $(x^{-1},0)=(1,0)\cdot(x^{-1},y^{-1})$ なども像に入る．$(1,0),(0,1),(x^{\pm1},0),(0,y^{\pm1})$ は右辺を環として生成するので全射．
  \item \emph{単射．}$\tilde\phi(h/(x+y)^m)=\phi(h)\cdot(\text{単元})$ なので，$\tilde\phi(h/(x+y)^m)=0$ なら $h(x,0)=0$ かつ $h(0,y)=0$．
        $h$ の代表元 $H\in k[x,y]$ は $y$ でも $x$ でも割り切れるので $H\in(xy)$，すなわち $R$ で $h=0$．
\end{itemize}
よって $\tilde\phi$ は同型であり，$e$ による分解 $R_{x+y}=R_{x+y}e\times R_{x+y}(1-e)$ がこの直積分解に対応する．

\textbf{幾何学的な比較．}
$\Zero(xy)$ 上で $x+y$ は，$x$ 軸 $\{(a,0)\}$ 上では $a$，$y$ 軸 $\{(0,b)\}$ 上では $b$ に等しい．したがって $x+y$ が消えるのは原点だけであり，
\[
  D(x+y)=\Zero(xy)\setminus\{(0,0)\}=\bigl(x\text{ 軸}\setminus\{0\}\bigr)\ \sqcup\ \bigl(y\text{ 軸}\setminus\{0\}\bigr)
\]
は，交わらない二つの $\A^1\setminus\{0\}$ の非交和である（図\ref{fig:ch06-sol-Dxy}）．
各成分の関数環は $k[x,x^{-1}]$，$k[y,y^{-1}]$（第5章の例，$\A^1\setminus\{0\}\cong\Zero(xy-1)$）であり，非交和の関数環は各成分の関数環の直積になる．
原点を除いたことで二本の軸が切り離され，その結果が環の直積分解（非自明な冪等元 $e$ の出現）として現れている．
例\ref{ex:ch06-loc-ex}(4) の $R_x\cong k[x,x^{-1}]$（$D(x)$ は $x$ 軸の部分だけ）と比べると，$R_{x+y}$ は両方の軸の部分を見ている．

\begin{figure}[H]
\centering
\begin{tikzpicture}[scale=0.8]
  \draw[very thick,blue] (-2.2,0)--(-0.12,0);
  \draw[very thick,blue] (0.12,0)--(2.2,0) node[right]{$x$ 軸 $\setminus\{0\}$：$k[x,x^{-1}]$};
  \draw[very thick,red] (0,-1.8)--(0,-0.12);
  \draw[very thick,red] (0,0.12)--(0,1.8) node[above]{$y$ 軸 $\setminus\{0\}$：$k[y,y^{-1}]$};
  \draw[thick,fill=white] (0,0) circle (3pt);
\end{tikzpicture}
\caption{$D(x+y)\subset\Zero(xy)$：原点を除くと二つの成分に分かれる．}
\label{fig:ch06-sol-Dxy}
\end{figure}

\exsol{prob:ch06-2-2}
加法は二つの引数について対称な式で定義されているので，第1引数の代表元を取り替えても値が変わらないことを示せば十分である（第2引数については対称性から同じ議論が適用でき，両方を取り替える場合は一つずつ取り替えればよい）．

$r/s=r_1/s_1$，すなわちある $u\in S$ で $u(rs_1-r_1s)=0$ とする．示すべきは
\[
  \frac{rs'+r's}{ss'}=\frac{r_1s'+r's_1}{s_1s'}
\]
であり，そのためには $S$ のある元を掛けて $(rs'+r's)s_1s'-(r_1s'+r's_1)ss'$ が $0$ になればよい．展開すると
\[
  (rs'+r's)s_1s'-(r_1s'+r's_1)ss'
  =rs_1s'^2+r'ss_1s'-r_1ss'^2-r'ss_1s'
  =s'^2\,(rs_1-r_1s)
\]
（$r'$ を含む項は打ち消し合う）．これに $u\in S$ を掛けると $s'^2\cdot u(rs_1-r_1s)=0$．
$u\in S$ なので定義により二つの分数は等しい．以上で和が代表元のとり方によらないことが示された．
（ここでも，$rs_1-r_1s$ 自体は $0$ とは限らず，$u$ を掛けて初めて $0$ になることが本質的である．）

\exsol{prob:ch06-3-1}
$f=y^2-x^2(x+1)$ は既約多項式（例\ref{ex:ch06-node-branch}）なので $R=k[x,y]/(f)$ は整域である．
$q=(-1,0)$，$p=(0,0)$ はともに $\Zero(f)$ の点であり，$\m_q=(x+1,\ y)$，$\m_p=(x,\ y)$（$R$ のイデアルとして）．

\textbf{$R_{\m_q}$ は離散付値環．}
$R$ は Noether 環なのでその局所化 $R_{\m_q}$ も Noether（$R_{\m_q}$ のイデアルは $R$ のイデアルの拡大 $S^{-1}\af$ であり，$\af$ の有限生成系がそのまま生成系になる）であり，整域 $R$ の局所化なので局所整域である（命題\ref{prop:ch06-Rp-local}）．
極大イデアルは $\m_qR_{\m_q}=(x+1,\ y)$ である．$R$ において
\[
  y^2=x^2(x+1)
\]
であり，$x(q)=-1\ne0$ なので $x\notin\m_q$，すなわち $x$ は $R_{\m_q}$ の単元である．よって
\[
  x+1=\frac{y^2}{x^2}\in(y)R_{\m_q},
\]
したがって $\m_qR_{\m_q}=(y)$ は単項である．また $R$ は整域で $y\ne0$（$y\notin(f)$）なので $y\ne0$ in $R_{\m_q}$ であり，極大イデアルは $0$ でない．
以上より $R_{\m_q}$ は離散付値環である．$y$ が局所パラメータであり，例えば $\mathrm{ord}_q(x+1)=2$（$x+1=y^2\cdot x^{-2}$）：直線 $x=-1$ は曲線に $q$ で接している．
（幾何学的には，$q$ の近くでは $x\ne0$ なので $y\mapsto$（曲線上の点）が局所座標になり，$q$ は曲線の非特異点である．）

\textbf{$R_{\m_p}$ の極大イデアルは単項でない．}
$A=R_{\m_p}$，$\m=\m_pA=(x,y)$，$\kappa=A/\m\cong k$ とする．$\dim_k\m/\m^2=2$ を示す．
$\m$ は $x,y$ で生成されるので，$\m/\m^2$ は $\bar x,\bar y$ で張られる．一次独立性を示すために $ax+by\in\m^2$（$a,b\in k$）と仮定する．
$\m^2=\m_p^2A$ なので $ax+by=h/s$（$h\in\m_p^2$，$s\in R\setminus\m_p$）と書け，$R$ は整域なので $R$ において $s(ax+by)=h\in\m_p^2$．
これを $k[x,y]$ に持ち上げると，代表元 $S$（$S(0,0)=c\ne0$）について
\[
  S\cdot(ax+by)\in(x,y)^2+(f)=(x,y)^2
\]
（$f=y^2-x^2-x^3$ は2次以上の項しかもたないので $f\in(x,y)^2$）．左辺の1次の部分は $c(ax+by)$ であり，$(x,y)^2$ の元は1次の項をもたないので $a=b=0$．
よって $\dim_k\m/\m^2=2$．

もし $\m=(g)$ が単項なら $\m/\m^2$ は $\bar g$ で張られ，$\dim_k\m/\m^2\le1$ となって矛盾する．よって $\m_pR_{\m_p}$ は単項でなく，$R_{\m_p}$ は離散付値環でない．
曲線（1次元）なのに余接空間が2次元であることは，原点が結節点（特異点）であることを表す（系\ref{cor:ch06-min-gen}より $\m$ の極小生成系はちょうど2個）．
これは $q$ では $\m_q/\m_q^2$ が $\bar y$ で張られる1次元空間であることと対照的である．

\exsol{prob:ch06-4-1}
\textbf{$k[[x]]$ の場合．}
$0\ne f=\sum_{n}a_nx^n\in k[[x]]$ に対し，$a_n\ne0$ なる最小の $n$ を $\mathrm{ord}(f)$ とする．$f=x^n u$，$u=\sum_{m\ge0}a_{n+m}x^m$ と書けば，$u$ の定数項 $a_n\ne0$ なので $u$ は単元（例\ref{ex:ch06-local-ex}(4)）．
$I\ne0$ をイデアルとし，$n=\min\{\mathrm{ord}(f)\mid0\ne f\in I\}$ とする．最小を与える $f=x^nu\in I$ について $x^n=fu^{-1}\in I$ なので $(x^n)\subset I$．
逆に $0\ne g\in I$ は $\mathrm{ord}(g)=m\ge n$ なので $g=x^mv=x^n\cdot x^{m-n}v\in(x^n)$．よって $I=(x^n)$．
したがって $k[[x]]$ のイデアルは $0$ と $(x^n)$（$n\ge0$，$(x^0)=k[[x]]$）に限る．また $\mathrm{ord}$ を比べれば $(x^n)$ たちは互いに異なり，$k[[x]]\supsetneq(x)\supsetneq(x^2)\supsetneq\cdots$ となる．

\textbf{$k[x]_{(x)}$ の場合．}
$0\ne\varphi=f/g\in k[x]_{(x)}$（$g(0)\ne0$）について，$f=x^nf_1$（$f_1(0)\ne0$）と書くと
\[
  \varphi=x^n\cdot\frac{f_1}{g},\qquad \frac{f_1}{g}\text{ は単元（逆元 }g/f_1\in k[x]_{(x)}\text{）}.
\]
$n$ は $\varphi$ の原点での零点の位数 $\mathrm{ord}_0(\varphi)$ である（表示 $x^nu$ の $n$ は一意：$x^nu=x^mv$，$n<m$ なら $u=x^{m-n}v$ が単元でなくなり矛盾）．
以下，$k[[x]]$ の場合と全く同じ議論で，イデアル $I\ne0$ は $n=\min\mathrm{ord}_0$ に対し $I=(x^n)$ となる．
（あるいは，問\ref{prob:ch06-1-1}の単射 $T\colon k[x]_{(x)}\to k[[x]]$ が $\mathrm{ord}$ を保つことに注意してもよい．）
これらの環はいずれも離散付値環であり，$\mathrm{ord}$ がその付値である．

\exsol{prob:ch06-4-2}
$e^2=e$ とすると $e(1-e)=0$．また $e+(1-e)=1\notin\m$ なので，$e$ と $1-e$ の少なくとも一方は $\m$ に属さず，局所環では $\m$ の外の元は単元である（命題\ref{prop:ch06-local-char}）．
$e$ が単元なら $e(1-e)=0$ の両辺に $e^{-1}$ を掛けて $1-e=0$，すなわち $e=1$．$1-e$ が単元なら同様に $e=0$．

\textbf{非連結な空間への応用．}
空間 $Y$ が空でない開かつ閉な部分集合 $Y_1,Y_2$ の非交和 $Y=Y_1\sqcup Y_2$ であるとき，$Y_1$ 上で $1$，$Y_2$ 上で $0$ をとる関数 $e$ は局所的に定数なので正則関数であり，
$e^2=e$，$e\ne0,1$ を満たす．したがって $Y$ の正則関数環は非自明な冪等元をもち，上の結果により局所環ではありえない．

$Y=\Zero(xy)\setminus\{0\}$ の場合，問\ref{prob:ch06-2-1}で見たように $Y=D(x+y)$ であり，その関数環は
$R_{x+y}\cong k[x,x^{-1}]\times k[y,y^{-1}]$ である．冪等元 $e=x/(x+y)\leftrightarrow(1,0)$ が非自明なので局所環でない．
実際この環は少なくとも二つの極大イデアル（例えば $\{0\}\times k[y^{\pm1}]+k[x^{\pm1}]\cdot(x-1)$ に相当する点 $(1,0)$ と点 $(0,1)$ の極大イデアル）をもつ．
幾何学的には，局所環は「一点のまわり」の関数環なので\emph{局所的に連結}であり，二つ以上の連結成分に分かれた空間を表すことはできない．
一方 $\Zero(xy)$ 自身の原点での局所環 $R_{\m_0}$ は（原点で二本の軸がつながっているので）冪等元 $0,1$ しかもたず，局所環である（例\ref{ex:ch06-node-local}）．

\exsol{prob:ch06-5-1}
$\OO=\OO_{\A^2,0}=k[x,y]_{(x,y)}$，$\m=(x,y)\OO$ とする．

\textbf{準備：多項式が $\m^n$ に属する条件．}
多項式 $h$ について，$h\in\m^n$ $\iff$ $h$ の最低次の斉次成分の次数が $n$ 以上．
実際 $\Leftarrow$ は $(x,y)^n$ が $n$ 次以上の多項式全体であることから明らか．$\Rightarrow$：$h\in\m^n=(x,y)^n\OO$ なら $h=a/s$（$a\in(x,y)^n$，$s(0,0)\ne0$）と書け，
$k[x,y]$ は整域なので $sh=a\in(x,y)^n$．$sh$ の最低次の成分は $s(0,0)\cdot(h\text{ の最低次の成分})$ なので，$h$ の最低次は $n$ 以上．

\textbf{$n$ の決定．}
$f=x^2+y^3$ の最低次の成分は $x^2$（2次）なので $f\in\m^2\setminus\m^3$，$n=2$．
$g=x+y^2$ の最低次の成分は $x$（1次）なので $g\in\m\setminus\m^2$，$n=1$．

\textbf{$\m/\m^2$ の次元．}
$\OO/(h)$ の極大イデアルは $\bar\m=\m/(h)$ であり，$\bar\m^2=(\m^2+(h))/(h)$ なので
\[
  \bar\m/\bar\m^2\cong\m/(\m^2+(h)).
\]
上の準備から $\m/\m^2$ は $\bar x,\bar y$ を基底とする2次元空間である．
\begin{itemize}
  \item $h=f$：$f\in\m^2$ なので $\m^2+(f)=\m^2$ であり，$\dim_k\bar\m/\bar\m^2=2$．
  \item $h=g$：$g\equiv x\pmod{\m^2}$ なので $\m/(\m^2+(g))=\m/(\m^2+(x))$ は $\bar y$ で張られ，$\dim_k\bar\m/\bar\m^2=1$．
\end{itemize}

\textbf{特異性の判定．}
$\OO/(h)\cong\OO_{\Zero(h),0}$（局所化と剰余の可換性）であり，$f,g$ はともに既約（$x$ についてモニック2次式・1次式とみれば，$f$ が分解するには $a(y)^2=-y^3$ が必要で不可能，$g$ は $x$ について1次）なので $\Zero(f)$，$\Zero(g)$ は曲線（1次元）である．
曲線上の点が非特異であるのは余接空間 $\bar\m/\bar\m^2$ の次元が $1$ のときである（第7章）．よって
\begin{itemize}
  \item $\Zero(g)$（放物線 $x=-y^2$）は原点で非特異．実際 $\Zero(g)\cong\A^1$（$y\mapsto(-y^2,y)$）であり，$\bar\m=(\bar y)$ は単項で $\OO_{\Zero(g),0}$ は DVR である．
  \item $\Zero(f)$（$x^2=-y^3$，尖点曲線）は原点で特異．余接空間が2次元である．
\end{itemize}
一般に，原点を通る曲線 $\Zero(h)$ が原点で特異であることと $h\in\m^2$（$h$ の1次の部分が $0$，すなわち $\partial_xh(0)=\partial_yh(0)=0$）であることは同値であり，これが第7章の Jacobian 判定法の最も簡単な場合である．

\exsol{prob:ch06-6-1}
\textbf{$C^\infty_0$ は局所環．}
$\m=\{\varphi\in C^\infty_0\mid\varphi(0)=0\}$ は評価写像 $C^\infty_0\to\R$，$\varphi\mapsto\varphi(0)$（芽の値は well-defined で，定数関数があるので全射）の核なので，極大イデアルである．
$\varphi(0)\ne0$ なら，連続性により $\varphi$ は $0$ のある近傍 $U$ で $0$ にならず，$1/\varphi$ は $U$ 上 $C^\infty$ なので，その芽が $\varphi$ の逆元である．
よって $\m$ の外の元はすべて単元であり，命題\ref{prop:ch06-local-char}(3) により $(C^\infty_0,\m)$ は局所環である．

\textbf{$\m^n$ の記述（Hadamard の補題）．}
$\varphi$ が $\varphi(0)=\varphi'(0)=\cdots=\varphi^{(n-1)}(0)=0$ を満たせば，積分形の剰余項をもつ Taylor の公式により
\[
  \varphi(x)=\frac{x^n}{(n-1)!}\int_0^1(1-s)^{n-1}\varphi^{(n)}(sx)\,ds=x^n\varphi_n(x)
\]
と書け，$\varphi_n$ は $C^\infty$ である（積分記号下の微分）．特に $n=1$ で $\m=(x)$ であり，$\m^n=(x^n)$．
逆に $x^n\psi$ は $n-1$ 階までの微分が原点で $0$ なので，
\[
  \m^n=\{\varphi\in C^\infty_0\mid\varphi^{(j)}(0)=0\ (0\le j<n)\},\qquad
  \m^\infty=\bigcap_n\m^n=\{\varphi\mid\text{原点での全ての微分係数が }0\}.
\]
したがって $C^\infty_0/\m^\infty$ への像は，$\varphi$ の Taylor 級数 $\sum\varphi^{(j)}(0)x^j/j!\in\R[[x]]$ と同じ情報である．

\textbf{$e^{-1/x^2}$ の芽．}
$\psi(x)=e^{-1/x^2}$（$x\ne0$），$\psi(0)=0$ とする．$x\ne0$ では帰納法により $\psi^{(n)}(x)=P_n(1/x)\,e^{-1/x^2}$（$P_n$ は多項式）と書ける：
$\frac{d}{dx}\bigl(P(1/x)e^{-1/x^2}\bigr)=\bigl(-P'(1/x)x^{-2}+2P(1/x)x^{-3}\bigr)e^{-1/x^2}$．
また任意の $m\ge0$ について $t=1/x^2$ とおけば $|x|^{-m}e^{-1/x^2}=t^{m/2}e^{-t}\to0$（$x\to0$）．
よって帰納的に $\psi^{(n)}(0)=\lim_{x\to0}\psi^{(n-1)}(x)/x=\lim_{x\to0}(x^{-1}P_{n-1}(1/x))e^{-1/x^2}=0$ となり，$\psi$ は $C^\infty$ で全ての微分係数が $0$．
すなわち $\psi\in\m^\infty$ であり，Taylor 展開は $0$ である．一方 $x\ne0$ では $\psi(x)>0$ なので，原点のどんな近傍でも $\psi$ は恒等的に $0$ でなく，芽 $\psi\ne0$．
したがって $C^\infty_0$ では $\m^\infty\ne0$ であり，Taylor 展開 $C^\infty_0\to\R[[x]]$ は単射でない．

\textbf{$k[x]_{(x)}$ との違い．}
$k[x]_{(x)}$ では問\ref{prob:ch06-4-1}により $\m=(x)$，$\m^n=(x^n)$ であり，零でない元は $x^nu$（$u$ 単元，$n$ は一意）と書けるので $x^{n+1}$ の倍数ではない．
よって $\bigcap_n\m^n=0$ であり，Taylor 展開 $T\colon k[x]_{(x)}\to k[[x]]$（問\ref{prob:ch06-1-1}）は単射である：\emph{代数的な芽は Taylor 級数で完全に決まる．}
一般に Noether 局所環 $(R,\m)$ では $\bigcap_n\m^n=0$（Krull の交叉定理）が成り立つ．
$C^\infty_0$ で $\m^\infty\ne0$ となることは，$C^\infty_0$ が Noether 環でないことを示している（Krull の交叉定理の対偶）．

この違いは本章の「関数の剛性」の議論（\ref{sec:ch06-phys}節 (b)）そのものである．
$C^\infty$ 関数は「Taylor 展開では見えない」平坦な関数 $\psi$ をもち，これを使って bump function が作れる（台の外では全ての微分が消える）．
代数的な関数にはそのような自由度がなく，$\m$ 進的に見て原点の無限小近傍が関数を決めてしまう．
物理では，結合定数 $g$ についての摂動級数がすべての次数で $0$ になる非摂動効果 $e^{-1/g^2}$（インスタントン効果など）がまさに $\psi$ と同じ型であり，
摂動展開（$\m$ 進展開）からは原理的に見えない寄与が存在することの数学的な表現になっている．
